\section{Practical checks}
\label{sec:validation}

Before increasing the size of a calculation, check these four choices:
\begin{description}
\item[Units.] Use one time unit throughout. Frequencies, decay rates and
\texttt{eta} are in its inverse; pulse widths and delays are in that unit.
With physical Hamiltonian energies, use the matching value of $\hbar$.
\item[Initial state.] Pass the intended density matrix explicitly.
For stationary preparation, use \texttt{rho0=rho\_ss} after computing
\texttt{rho\_ss = system.stationary\_state()}.
\item[Frequency grid.] Start with a small grid, then reduce the spacing
around peaks until their shape and amplitude stop changing appreciably.
The result is indexed as \texttt{S[omega3\_index, omega1\_index]}.
\item[Convergence parameter.] Use a small positive \texttt{eta} to ensure
causal-transform convergence. Reduce it while refining the frequency grid
until its added broadening no longer materially changes the features of
interest and the linear solves remain stable. Check unnormalized complex
responses so normalization does not hide amplitude changes.
\end{description}

Detailed checks of the implementation and time--frequency comparisons
are kept separately in \path{docs/VALIDATION.md} and
\path{docs/advanced_notes.md}.
